\usection{Justificación}

La era digital actual se define por la rápida evolución de la tecnología, lo que demanda que los estudiantes no sólo sean consumidores de tecnología, sino también creadores y pensadores críticos capaces de comprender y moldear el mundo que los rodea. Para lograr esto, es fundamental desarrollar el pensamiento computacional: un conjunto de habilidades para la resolución de problemas que incluye la descomposición, el reconocimiento de patrones, la abstracción y el diseño de algoritmos.

Este proyecto se distingue por su originalidad al proponer el desarrollo de un lenguaje de programación visual basado en el dataflow, diseñado específicamente para la enseñanza del pensamiento computacional. A diferencia de los enfoques tradicionales basados en texto, un lenguaje visual permite a los estudiantes construir programas arrastrando y conectando bloques funcionales, lo que proporciona una representación directa y clara de la lógica del programa.

Ofrecer una alternativa a los lenguajes visuales basados en flujo de control es importante para adaptarse a las habilidades, capacidades y preferencias de los estudiantes, brindando distintos medios para la formación de una cualidad tan importante en la actualidad. Esto se alinea con las ideas de \citeA{meyer2014} que plantean un Diseño Universal para el Aprendizaje, proponiendo el uso de esquemas flexibles para garantizar oportunidades de aprendizaje equitativas. Este se sustenta en parte en un principio de múltiples medios de representación, lo que apoya la coexistencia de distintos medios de programación visual.

Además, \citeA{wadge1985} defienden que los lenguajes basados en dataflow ayudan a visualizar de forma sencilla la interacción entre módulos, a la vez que elimina el vacío entre el problema y la solución, cosa que es valiosa en la formación del pensamiento computacional. Igualmente, este tipo de lenguajes facilita la comprensión de conceptos avanzados como la concurrencia y los sistemas distribuidos, que son cada vez más relevantes en la computación moderna, y que en el contexto educativo actual no son completamente soportados por herramientas basadas en flujo de control.

Similar a como indica \citeA{resnicketal2009} sobre el diseño de Scratch, el propósito de este proyecto es empoderar a los jóvenes a pensar creativamente, a razonar sobre problemas y a trabajar colaborativamente; con la meta de formar el pensamiento computacional. Es importante recordar que aprender a programar no es la meta fundamental por sí misma, sino que es la vía por la cual se desarrolla un modo de pensamiento que permite resolver problemas y desenvolverse en una sociedad de información donde las tecnologías y conceptos de computación están siempre presentes.

Por otro lado, el lenguaje desarrollado con este propósito podría tener aplicaciones generales integradas con tendencias tecnológicas actuales, como la implementación de agentes de IA, la programación de robots y la codificación de entornos de Internet de las Cosas (IoT). Adicionalmente, este proyecto busca fomentar la investigación sobre nuevas metodologías pedagógicas en la educación en programación, lo que podría motivar el desarrollo de herramientas similares enfocadas en otros conceptos avanzados de la informática.

Este proyecto contribuiría directamente al Objetivo de Desarrollo Sostenible (ODS) 4: Educación de Calidad, específicamente a la meta 4.4, que busca aumentar considerablemente el número de jóvenes y adultos que tienen las competencias necesarias, en particular técnicas y profesionales, para acceder al empleo, el trabajo decente y el emprendimiento. Al facilitar el desarrollo del pensamiento computacional y las habilidades de programación, se prepararía a los estudiantes con competencias clave requeridas en el siglo XXI. Adicionalmente, los educadores son otro grupo beneficiado, ya que contarían con un medio didáctico diseñado específicamente para enseñar conceptos de computación de manera más accesible y atractiva para los estudiantes.

Además, al crear una herramienta educativa innovadora y accesible, se estaría apoyando al ODS 9: Industria, Innovación e Infraestructura, al promover la innovación tecnológica en el sector educativo y fomentar una cultura de desarrollo y creatividad desde una edad temprana. En este ODS también se aprecia como tarea fundamental la reducción de la brecha digital, para lo que la educación en computación resulta de suma importancia.